\documentclass[10pt]{article}

\usepackage[utf8]{inputenc}

\usepackage[T1]{fontenc}

\usepackage{amsmath}

\usepackage{amsfonts}

\usepackage{amssymb}

\usepackage[version=4]{mhchem}

\usepackage{stmaryrd}

\usepackage{bbold}

\usepackage{graphicx}

\usepackage[export]{adjustbox}

\graphicspath{ {./images/} }



\begin{document}

нений макроскопич. физики. П.о. к. П., в частности, удовлетворяет равновесный статистич. оператор Гиббса.



Лит.: [1] Б о го л ю б о в Н. Н., Избранные труды, т. 3, Киев, 1971; [2] Боголюбов Н. Н., Боголюбов Н.Н. (мл.), Введение в квантовую статистическую механику, М., 1984. М.Ю. Ковалевский.



ПРОСТРАНСТВЕННОПОДОБНЫЙ BEКТОР в ч ас тной (специальной) и общей теории относи тельности - четырехмерный вектор, сумма квадратов пространственных компонент которого больше квадрата его временной компоненты. П. в., имеющий начало в какой-нибудь точке четырехмерного пространства-времени, лежит вне внутренних полостей светового конуса с вершиной в данной точке. Всегда существует система отсчета, в к-рой временная компонента П. в. обращается в нуль, и у него остаются только пространственные компоненты. В Минковского пространстве-времени с метрич. тензором $(+1,-1,-1,-1)$ квадрат длины П. в. $A$ отрицателен:



\[

(A)^{2}=A^{\mu} A_{\mu}=\left(A^{0}\right)^{2}-A^{2}<0, \mu=0,1,2,3 .

\]



Здесь $A^{0}$ - временная, $A^{i}, i=1,2,3,-$ пространственная компоненты 4-вектора, $A=\left(A^{1}, A^{2}, A^{3}\right)$. ПРОСТРАНСТВО КОНФИГУРАЦИЙ - см. Классический ансамбль.



ПРОСТРАНСТВО РЕАЛИЗАЦИЙ - см. СлучайнЫй процесс. ПРОСТРАНСТВО С ИНДЕФИНИТНОЙ МЕТРИКОЙ векторное пространство над полем $\mathbb{C}$ (соответственно $\mathbb{R})$ c полуторалинейной (соответственно билинейной) формой, которая, не являясь положительно определенной, по традиции называется и нде ф и н и т н ой ме т р и кой. Исследование геометрии развито для П.с и. м., имеющих добавочную структуру топологич. пространства.



$\Pi$ р и мер. Так называемое пространство с $G$-мет ри к о й - гильбертово пространство $\boldsymbol{H}$ с эрмитовой непрерывной формой $\langle x, y\rangle=(x, G y)$, где $(\cdot, \cdot)-$ скалярное произведение, а оператор $\boldsymbol{G}$ ограничен и самосопряжен. Наиболее продвинута теория операторов в $J$-пространствах, когда индефинитная метрика в гильбертовом пространстве $\langle x, y\rangle=(x, J y)$ задается такой инволюцией $J=J^{-1}=J^{*}$, что $J=P_{+}-P_{-}$; здесь $P_{+}-$ортогональные проекторы, $P_{ \pm} \neq 0$, $P_{+} P_{-}=0, P_{+}+P_{-}=I$ и $\langle x, y\rangle=\left(P_{+} x, P_{+} y\right)-\left(P_{-} x, P_{-} y\right)$. При



\[

x=\min \left\{\operatorname{dim} P_{+} H, \operatorname{dim} P_{-} H\right\}<\infty

\]



$J$-пространство называется пространством Понтрягина $\Pi_{x}$, при $x=\infty-$ пространством Крейна. Пространства Крейна применяются, в частности, в квантовой теории калибровочных полей.



Лит.: [1] Ги н з бург Ю. П., Иох в идо в И. С., «Успехи математич. наук», 1962, т. 17 , в. 4, с. 3-56; [2] Крей н М. Г., в Кн.: II летняя математическая школа, т. 1, Киев, 1965, с. 15-92; [3] Н а д ь К., Пространства состояний с индефинитной метрикой в квантовой теории поля, пер. с англ., М., 1969; [4] В оg n a r J., Indefinite inner product spaces, B.- [a. o.], 1974; [5] Аз изов Т.Я., Иохв идов И.С., Основы теории линейных операторов в пространствах с индефинитной метрикой, М., 1986; [6] Хоруж и й С. С., Введение в алгебраическую квантовую теорию поля, М., $1986 . \quad$ А. В. Булинский. ПРОСТРАНСТВО С УТЛОВОЙ МЕТРИКОЙ - см. Вейля связность.



IPOCTPAНСТВО TИПА $\boldsymbol{S}^{\prime}$ - пространство обобщенных функций, обладающих свойствами ультрараспределений и умеренных распределений. Введены И. М. Гельфандом и Г. Е. Шиловым в конце 50-х гт. 20 в.



Основное пространство $S_{p}, p \geqslant 0$, состоит из всех функций $\varphi(x)$, удовлетворяющих оценкам



\[

\left|x^{\alpha} D^{\beta} \varphi(x)\right| \leqslant C_{|\beta|^{A}}{ }^{|\alpha|}|\alpha|^{p|\alpha|},|\alpha|,|\beta|=0,1,2, \ldots, x \in \mathbb{R}^{n},

\]



где постоянные $C_{|\beta|}, A$ зависят от $\varphi$.



Основное пространство $S^{q}, q \geqslant 0$, состоит из всех функций $\varphi(x)$, удовлетворяющих оценкам



\begin{center}

\includegraphics[max width=\textwidth]{2023_07_08_16345ca3c613250d2b33g-1}

\end{center}



где постоянные $C_{|\alpha|}, B$ зависят от $\varphi$. Основное пространство $S_{p}^{q}, p, q \geqslant 0$, состоит из всех фунКций $\varphi(x)$, удовлетворяющих оценкам



\begin{center}

\includegraphics[max width=\textwidth]{2023_07_08_16345ca3c613250d2b33g-1(1)}

\end{center}



\[

\begin{aligned}

& |\alpha|,|\beta|=0,1,2, \ldots, x \in \mathbb{R}^{n},

\end{aligned}

\]



где постоянные $C, A, B$ зависят от $\varphi$.



Здесь функции $\varphi(x)$ определены и имеют непрерывные производные всех порядков в $\mathbb{R}^{n}, \alpha=\left(\alpha_{1}, \ldots, \alpha_{n}\right), \beta=\left(\beta_{1}, \ldots, \beta_{n}\right)$,мультииндексы, $x^{\alpha}=x_{1}^{\alpha_{1}} \ldots x_{n}^{\alpha_{n}}, D^{\beta}=D_{1}^{\beta_{1}} \ldots D_{n}^{\beta_{n}}, D_{j}=\partial / \partial x_{j}$, $|\alpha|=\alpha_{1}+\ldots+\alpha_{n},|\beta|=\beta_{1}+\ldots+\beta_{n}$.



Пространства $S_{p}, S^{q}, S_{p}^{q}$ снабжаются естественными локально выпуклыми топологиями. Сопряженные пространства называются пространствами типа $S^{\prime}$. Пространства $S_{p}$ и $S^{q}$ всегда нетривиальны, то есть содержат отличные от нуля элементы, пространства $S_{p}^{0}$ и $S_{0}^{q}$ нетривиальны тогда и только тогда, когда $p>1$ и $q>1, S_{p}^{q}$ при $p>0$ и $q>0$ нетривиально тогда и только тогда, когда $p+q \geqslant 1$. Преобразование Фурье $F$ (см. Фурье преобразование обобщенных функций) задает топологич. изоморфизмы



\[

F: S_{p} \cong S^{p}, F: S_{p}^{q} \cong S_{q}^{p}

\]



и аналогичные изоморфизмы для сопряженных пространств. В пространствах $S^{q}$ и $S_{p}^{q}, q>0$, действуют дифференциальные операторы бесконечного порядка



\[

P(D)=\sum_{|\alpha|=0}^{\infty} a_{\alpha} D^{\alpha},

\]



где символ $P(z)$ - целая функция минимального типа порядка роста $1 / q$.



Изучаются пространства типа $S^{\prime}$ более общего вида с заменой последовательности $|\alpha|^{p|\alpha|}$ на $a_{|\alpha|}$ и $|\beta|^{q|\beta|}$ на $b_{|\beta|}$.



Лит:: [1] Гельфанд И. М., Шилов Г. Е., Пространства основных и обобщенных функций, М., 1958.



В. В. Жаринов.



ПРОСТPАНСТВО-ВРЕМЯ Ньютоновской Грави т а ц и и - геометрическая модель гравитационного поля, представляющая собой прямое произведение $E^{1} \times E^{3}$ оси времени $E^{1}$ на евклидово пространство $E^{3}$. Точки П.-в. ньютоновской гравитации моделируют события. Расстояние между двумя точками в П.-в. ньютоновской гравитации определено лишь для одновременных событий. П.-в. ньютоновской гравитации, таким образом, представляет собой расслоенное по времени 4-мерное пространство-время со слоями $E^{3}$. Свободное движение (движение по инерции) в П.-в. ньютоновской гравитации происходит по кривым, соответствующим свободному движению в гравитационном поле. Различным ньютоновским гравитационным полям соответствуют различные П.-в. ньютоновской гравитации. Величина гравитационного поля характеризуется отклонением «геодезических» в соответствующем П.-в. ньютоновской гравитации от прямых, представляющих собой геодезические в пространстве-времени Галилея, к-рое является частным случаем П.-в. ньютоновской гравитации при отсутствии гравитационного поля и обладает десятипараметрич. группой движений Галилея. Световой конус в П.-в. ньютоновской гравитации вырождается, и в силу этого возможно движение с любыми скоростями.



Лит.: [1] Пе н р о у 3 Р., Структура пространства-времени, пер. с англ., М., $1972 . \quad$ А.Т. Ильичев. ПРОСТРАНСТВО-BPEMЯ те ории относительности - модель физического мира, гладкое четырехмерное многообразие с гладкой (псевдо)метрикой - интервалом. В основе концепции П.-в. лежит понятие то чеч н о го с о быт ия. Точечное событие моделируется точкой П.-в. и





\end{document}